\chapter{Quantum Computing: Amplitudes, Interference, and What Qubits Actually Do}
\label{ch42:quantum}

A popular myth pervades science journalism: ``A classical bit is either $0$ or
$1$, whereas a quantum bit can be $0$ and $1$ at the same time, so an $n$-qubit
computer tests all $2^n$ possible solutions in parallel simultaneously!''

This slogan is completely misleading. It sounds exhilarating, but it conceals
the most important physical rule of quantum mechanics: \textbf{measurement}.
Suppose you prepare an $n$-qubit quantum register in an equal blend of all
$2^n$ possible $n$-bit strings, and you measure it. You do not see all $2^n$
answers. Nature does not print a list of $2^n$ results on your screen. You
obtain \textbf{exactly one single bitstring}, chosen completely at random with
probability $1/2^n$.

If all a quantum computer could do was prepare a massive superposition and pick a
random item upon measurement, it would be nothing more than an exorbitantly
expensive physical dice roller. You could simulate that on your laptop in two lines
of code using a pseudo-random number generator, with zero quantum hardware.

What, then, does a quantum computer actually do? The secret lies in a profound
mathematical difference between classical probabilities and quantum amplitudes:
\begin{itemize}
  \item \textbf{Classical probabilities are non-negative numbers ($p \ge 0$).}
  If an algorithm has two distinct computational pathways that lead to the same
  incorrect output, the total probability of failing is $p_1 + p_2$. Because
  probabilities can never be negative, adding an alternative path to an error
  always increases (or at best maintains) your chance of failure.
  \item \textbf{Quantum states are described by probability amplitudes.}
  Amplitudes can be positive, negative, or complex numbers ($\alpha \in \C$). If
  a quantum algorithm provides two computational pathways that lead to the
  \emph{same incorrect output}, but one path arrives with amplitude $+a$ and the
  other path arrives with amplitude $-a$, the two amplitudes add linearly:
  \[
    +a + (-a) = 0.
  \]
  The two paths \textbf{destructively interfere and cancel each other out to exact
  zero}.
\end{itemize}

Quantum algorithm design is the art of choreographing a linear-algebraic dance
in which computational pathways leading to wrong answers interfere destructively
and cancel, while pathways leading to the correct answer interfere
constructively and reinforce. When you measure the machine at the end, the
surviving amplitude is concentrated almost entirely on the solution.

\begin{historybox}
The realization that quantum mechanics might possess unprecedented computational
power emerged from a frustrating roadblock in physics. In 1981, at a joint MIT--IBM
conference on the physics of computation, Richard Feynman pointed out that
simulating a system of $n$ interacting quantum particles on a classical digital
computer appears to require exponential resources. Because each quantum particle
adds a degree of freedom, describing the state of $n$ entangled particles requires
tracking $2^n$ continuous complex numbers. Simulating just a few hundred electrons
demands more memory than there are atoms in the observable universe. Feynman proposed
turning this curse into an opportunity: instead of struggling to simulate quantum
physics on classical logic gates, build a computing device whose physical elements
obey quantum laws, allowing nature to carry out the exponential linear algebra.
Concurrently, Paul Benioff proved that a quantum mechanical Hamiltonian could
simulate a classical reversible Turing machine.

In 1985, David Deutsch at Oxford laid the mathematical foundations of the field by
defining the Universal Quantum Computer, proving that a quantum machine could
evaluate functions in superposition. In 1992, Deutsch and Richard Jozsa formulated
the first quantum algorithm demonstrating a provable, deterministic separation
between quantum and classical query complexity. A year later, Ethan Bernstein and
Umesh Vazirani formalized the complexity class $\mathrm{BQP}$ (Bounded-Error Quantum
Polynomial Time) and presented an oracle problem establishing an exponential quantum
speedup under randomized query models.

For a decade, quantum computing remained largely a theoretical curiosity among
physicists and complexity theorists. That changed abruptly in 1994. Peter Shor,
working at AT\&T Bell Laboratories in New Jersey, announced polynomial-time quantum
algorithms for two famous mathematical problems: prime factorization and discrete
logarithms. Because the world's deployed public-key cryptographic infrastructure
(RSA, Diffie-Hellman, and elliptic curves) depended on the presumed intractability
of these two problems, Shor's algorithm transformed quantum computing from an
academic pursuit into an urgent global technological race. In 1996, Lov Grover,
also at Bell Labs, discovered an algorithm for searching an unsorted database of $N$
items in $\Theta(\sqrt{N})$ operations, establishing a quadratic speedup for
unstructured search that was soon proven by Bennett, Bernstein, Brassard, and
Vazirani to be the absolute mathematical limit for generic black-box search.

Early critics, including Rolf Landauer and William Unruh, raised severe physical
skepticism: quantum amplitudes are continuous numbers, so analog noise, stray
thermal photons, and environmental decoherence would inevitably scramble the delicate
phases before any non-trivial algorithm could complete. Unlike digital bits, which
snap back to clean voltage rails, analog states seemed doomed to continuous drift.
This skepticism was overcome when Peter Shor (1995) and Alexei Kitaev (1997)
invented quantum error correction and topological surface codes. They proved the
Quantum Threshold Theorem: if physical error rates per gate can be driven below a
constant threshold (roughly $1\%$), quantum errors can be actively detected and
corrected through discrete syndrome measurements without ever measuring or
collapsing the underlying quantum superposition. Today, laboratories worldwide
compete to build fault-tolerant hardware using superconducting transmon circuits,
trapped atomic ions, neutral atoms trapped in optical tweezers, and integrated
photonic waveguides.
\end{historybox}

\section{A Single Qubit as a 2D Unit Vector}
\label{ch42:single-qubit}

To understand quantum computing, you do not need quantum field theory, wave-particle
duality, or Hamiltonian operators. All you need is the linear algebra of unit
vectors in a two-dimensional vector space.

\subsection{Unwrapping Dirac's notation}

In classical computing, a bit stores a scalar value from the discrete set $\{0, 1\}$.
In quantum computing, the fundamental unit of information is a \textbf{qubit}
(quantum bit). A qubit lives in a two-dimensional vector space over the field of
complex numbers $\C^2$.

Physicists use Paul Dirac's notation, writing a vertical ket $\ket{\cdot}$ to denote
a column vector. The two standard basis vectors of $\R^2$ (or $\C^2$) are given
special labels:
\[
  \ket{0} = \colvec{1 \\ 0}, \qquad \ket{1} = \colvec{0 \\ 1}.
\]
Notice what this means: $\ket{0}$ is not the number zero; it is the unit vector pointing
along the horizontal $x$-axis. Likewise, $\ket{1}$ is the unit vector pointing along
the vertical $y$-axis. Any state of a single qubit is simply a linear combination of
these two orthogonal basis vectors:
\[
  \ket{\psi} = \alpha\ket{0} + \beta\ket{1} = \colvec{\alpha \\ \beta}.
\]
The coefficients $\alpha$ and $\beta$ are called \textbf{probability amplitudes}.

\subsection{Real amplitudes on the 2D unit circle}

Before we introduce complex numbers, let us build our intuition in the familiar
Cartesian plane $\R^2$, where $\alpha$ and $\beta$ are ordinary real numbers.

In $\R^2$, a qubit state $\ket{\psi} = \colvec{\alpha \\ \beta}$ is an arrow
originating at the origin. Nature imposes one fundamental constraint: the state
must be a \textbf{unit vector}. By the Pythagorean theorem, its Euclidean length
squared must equal $1$:
\[
  \norm{\ket{\psi}}^2 = \alpha^2 + \beta^2 = 1.
\]
A single qubit with real amplitudes is simply a point on the circumference of the
unit circle in the plane! For example, the state with equal positive amplitudes
is located at an angle of $45^\circ$:
\[
  \ket{+} = \frac{1}{\sqrt{2}}\ket{0} + \frac{1}{\sqrt{2}}\ket{1}
  = \colvec{\frac{1}{\sqrt{2}} \\ \frac{1}{\sqrt{2}}}.
\]
The state with opposite signs is located at an angle of $-45^\circ$ ($315^\circ$):
\[
  \ket{-} = \frac{1}{\sqrt{2}}\ket{0} - \frac{1}{\sqrt{2}}\ket{1}
  = \colvec{\frac{1}{\sqrt{2}} \\ -\frac{1}{\sqrt{2}}}.
\]
Both vectors have length $\sqrt{(1/\sqrt{2})^2 + (\pm 1/\sqrt{2})^2} = \sqrt{1/2 + 1/2} = 1$.

\subsection{Measurement and the Born Rule}

What happens when we observe a qubit? In classical mechanics, looking at an object
reveals its state without altering it. In quantum mechanics, observing a system
forces it to make a discrete, non-deterministic choice. This interaction is called
\textbf{measurement}.

\begin{definition}[The Born Rule in $\R^2$]
Let a qubit be in state $\ket{\psi} = \alpha\ket{0} + \beta\ket{1}$ with
$\alpha^2 + \beta^2 = 1$. When measured in the standard basis $\{\ket{0}, \ket{1}\}$:
\begin{enumerate}
  \item The outcome is $0$ with probability $\P(0) = \alpha^2$. If observed, the
  state collapses instantaneously onto $\ket{0} = \colvec{1\\0}$.
  \item The outcome is $1$ with probability $\P(1) = \beta^2$. If observed, the
  state collapses instantaneously onto $\ket{1} = \colvec{0\\1}$.
\end{enumerate}
\end{definition}

Geometrically, $\alpha$ is the horizontal projection (shadow) of the unit arrow onto
the $\ket{0}$-axis, and $\beta$ is the vertical projection (shadow) onto the
$\ket{1}$-axis. The probability of collapsing onto an axis is the \textbf{squared
length of that shadow}!

Because the state arrow has unit length, the Pythagorean theorem guarantees that the
probabilities sum to $1$:
\[
  \P(0) + \P(1) = \alpha^2 + \beta^2 = 1.
\]
If amplitudes are generalized to complex numbers $\alpha, \beta \in \C$, the exact
same geometry holds in $\C^2$: the probability of observing outcome $0$ is
$\abs{\alpha}^2$, the probability of observing outcome $1$ is $\abs{\beta}^2$, and
the normalization condition $\abs{\alpha}^2 + \abs{\beta}^2 = 1$ ensures that total
probability is conserved.

\begin{trapbox}
A common trap is confusing the probability amplitude $\alpha$ with the probability
$\abs{\alpha}^2$. If $\alpha = -1/\sqrt{2}$, the probability of measuring $0$ is
$(-1/\sqrt{2})^2 = +1/2$. The minus sign has zero effect on the immediate measurement
statistics of that single isolated qubit. But as we will see in the very next section,
that minus sign is stored in the state, and it causes massive computational
consequences when combined with subsequent operations.
\end{trapbox}

\section{Quantum Gates as Rotations and Reflections}
\label{ch42:quantum-gates}

How does a quantum computer perform a computation? In classical computing, logic
gates like \texttt{AND}, \texttt{OR}, and \texttt{NOT} transform input bits into output
bits. In quantum computing, a computation is performed by applying a linear operator
to the state vector.

\subsection{Conservation of probability and unitary matrices}

Every valid quantum gate must transform unit vectors into unit vectors. If a state
$\ket{\psi}$ satisfies $\norm{\ket{\psi}} = 1$, and we apply a matrix $U$, the resulting
vector $U\ket{\psi}$ must also satisfy $\norm{U\ket{\psi}} = 1$. Total probability
must remain $100\%$ under all circumstances.

In real vector spaces $\R^n$, matrices that preserve vector lengths are the
\textbf{orthogonal matrices} ($U\T U = I$). In complex vector spaces $\C^n$, matrices
that preserve vector lengths are the \textbf{unitary matrices}:
\[
  U^\dagger U = I,
\]
where $U^\dagger = (U\T)^*$ denotes the conjugate transpose (the adjoint).

This single geometric requirement has an immediate consequence: \textbf{every
isolated quantum computation is reversible}. Because $U^\dagger U = I$, the matrix
$U$ is always invertible, with its inverse given by its conjugate transpose:
$U^{-1} = U^\dagger$. A quantum gate never erases information.

\subsection{Elementary quantum gates}

Let us inspect the three foundational single-qubit gates:

\begin{enumerate}
  \item \textbf{The Pauli-$X$ Gate (Quantum NOT):}
  \[
    X = \colvec{0 & 1 \\ 1 & 0}.
  \]
  Applying $X$ to the basis states swaps them:
  \[
    X\ket{0} = \colvec{0 & 1 \\ 1 & 0}\colvec{1 \\ 0} = \colvec{0 \\ 1} = \ket{1},
    \qquad
    X\ket{1} = \colvec{0 & 1 \\ 1 & 0}\colvec{0 \\ 1} = \colvec{1 \\ 0} = \ket{0}.
  \]
  Geometrically, $X$ is a reflection across the diagonal line $y = x$.

  \item \textbf{The Pauli-$Z$ Gate (Phase Flip):}
  \[
    Z = \colvec{1 & 0 \\ 0 & -1}.
  \]
  Applying $Z$ leaves $\ket{0}$ unchanged and flips the sign of $\ket{1}$:
  \[
    Z\ket{0} = \ket{0}, \qquad Z\ket{1} = -\ket{1}.
  \]
  Geometrically, $Z$ is a reflection across the horizontal $\ket{0}$-axis.

  \item \textbf{The Hadamard Gate ($H$):}
  \[
    H = \frac{1}{\sqrt{2}}\colvec{1 & 1 \\ 1 & -1}.
  \]
  Notice that $H$ is symmetric ($H\T = H$) and orthogonal ($H\T H = I$), which means
  $H^2 = I$. The Hadamard gate is its own inverse!
\end{enumerate}

\subsection{Tracing the Hadamard gate by hand: The Two-Coin-Flip experiment}

The Hadamard gate is the engine of quantum superposition and interference. To see
why, let us walk through a thought experiment comparing a classical random process
with a quantum circuit.

Imagine a classical machine that flips a fair coin. We initialize the coin to Heads
($0$). The flipper scrambles it so that the state becomes Heads ($0$) with probability
$1/2$ and Tails ($1$) with probability $1/2$. Now we run the flipper a second time.
What is the outcome? A random coin flipped twice is still random:
$\P(0) = \frac{1}{2}\times\frac{1}{2} + \frac{1}{2}\times\frac{1}{2} = \frac{1}{4} + \frac{1}{4} = \frac{1}{2}$.
Randomness added to randomness produces randomness.

Now let us execute the exact same sequence on a quantum computer using the Hadamard
gate.

\textbf{Step 1: The first Hadamard gate.}
We initialize our qubit to $\ket{0} = \colvec{1\\0}$. We apply $H$:
\[
  H\ket{0} = \frac{1}{\sqrt{2}}\colvec{1 & 1 \\ 1 & -1}\colvec{1 \\ 0}
  = \colvec{\frac{1}{\sqrt{2}} \\ \frac{1}{\sqrt{2}}}
  = \frac{1}{\sqrt{2}}\ket{0} + \frac{1}{\sqrt{2}}\ket{1} = \ket{+}.
\]
If we choose to measure the qubit right now, what do we observe?
\begin{itemize}
  \item Outcome $0$ occurs with probability $(1/\sqrt{2})^2 = 1/2$.
  \item Outcome $1$ occurs with probability $(1/\sqrt{2})^2 = 1/2$.
\end{itemize}
The measurement statistics are identical to a fair coin flip!

What if we had initialized the qubit to $\ket{1} = \colvec{0\\1}$? Applying $H$ yields:
\[
  H\ket{1} = \frac{1}{\sqrt{2}}\colvec{1 & 1 \\ 1 & -1}\colvec{0 \\ 1}
  = \colvec{\frac{1}{\sqrt{2}} \\ -\frac{1}{\sqrt{2}}}
  = \frac{1}{\sqrt{2}}\ket{0} - \frac{1}{\sqrt{2}}\ket{1} = \ket{-}.
\]
If measured right now, the outcome is $0$ with probability $(1/\sqrt{2})^2 = 1/2$,
and $1$ with probability $(-1/\sqrt{2})^2 = 1/2$. To a direct measurement,
$\ket{+}$ and $\ket{-}$ look completely indistinguishable: both produce $50\%/50\%$
coin flips.

\textbf{Step 2: The second Hadamard gate without measurement.}
Now comes the demonstration of quantum interference. Suppose we start with
$\ket{0}$, apply $H$, do \textbf{not} measure, and immediately apply $H$ a second time.

Let us trace every individual computational path by distributing the linear operator:
\begin{align*}
  H(H\ket{0}) &= H\left(\frac{1}{\sqrt{2}}\ket{0} + \frac{1}{\sqrt{2}}\ket{1}\right) \\
  &= \frac{1}{\sqrt{2}}H\ket{0} + \frac{1}{\sqrt{2}}H\ket{1} \\
  &= \frac{1}{\sqrt{2}}\left(\frac{1}{\sqrt{2}}\ket{0} + \frac{1}{\sqrt{2}}\ket{1}\right)
   + \frac{1}{\sqrt{2}}\left(\frac{1}{\sqrt{2}}\ket{0} - \frac{1}{\sqrt{2}}\ket{1}\right) \\
  &= \left(\frac{1}{2}\ket{0} + \frac{1}{2}\ket{1}\right)
   + \left(\frac{1}{2}\ket{0} - \frac{1}{2}\ket{1}\right) \\
  &= \left(\frac{1}{2} + \frac{1}{2}\right)\ket{0} + \left(\frac{1}{2} - \frac{1}{2}\right)\ket{1} \\
  &= 1\ket{0} + 0\ket{1} = \ket{0}.
\end{align*}

Read that final line again. The amplitude of outcome $\ket{1}$ is
$\frac{1}{2} - \frac{1}{2} = 0$. The path through the intermediate state $\ket{0}$
contributed $+1/2$ to the final state $\ket{1}$, while the path through the
intermediate state $\ket{1}$ contributed $-1/2$ to the final state $\ket{1}$. The two
pathways destroyed each other completely!

Meanwhile, both intermediate pathways contributed $+1/2$ to the final state $\ket{0}$,
reinforcing each other to yield a total amplitude of $1$, giving a measurement
probability of $1^2 = 100\%$.

\begin{figure}[H]
\centering
\begin{tikzpicture}
  % Panel (a): Unit circle geometry
  \begin{scope}[shift={(0,0)}]
    \node[font=\small\bfseries] at (2.5, 6.1) {(a) Single Qubit on the 2D Unit Circle};
    
    % Axes
    \draw[->, gray!70, thick] (-0.3, 2.5) -- (5.2, 2.5) node[right, font=\footnotesize] {$\ket{0} = \colvec{1\\0}$};
    \draw[->, gray!70, thick] (2.5, -0.3) -- (2.5, 5.2) node[above, font=\footnotesize] {$\ket{1} = \colvec{0\\1}$};
    
    % Circle
    \draw[sourceblue!40, very thick] (2.5, 2.5) circle (2.2);
    
    % State vectors
    \draw[->, very thick, sourceblue!85!black] (2.5, 2.5) -- ++(45:2.2) node[above right, font=\footnotesize] {$\ket{+} = \tfrac{1}{\sqrt{2}}\ket{0} + \tfrac{1}{\sqrt{2}}\ket{1}$};
    \draw[->, very thick, quotientorange!90!black] (2.5, 2.5) -- ++(-45:2.2) node[below right, font=\footnotesize] {$\ket{-} = \tfrac{1}{\sqrt{2}}\ket{0} - \tfrac{1}{\sqrt{2}}\ket{1}$};
    \draw[->, very thick, gray!80!black] (2.5, 2.5) -- ++(0:2.2);
    \draw[->, very thick, gray!80!black] (2.5, 2.5) -- ++(90:2.2);
    
    % Angles
    \draw[dashed, gray!60] (2.5, 2.5) -- ++(22.5:2.55) node[right, font=\tiny] {Reflection axis ($22.5^\circ$)};
    \draw[sourceblue!80, thick] (3.3, 2.5) arc (0:45:0.8);
    \node[sourceblue!90!black, font=\scriptsize] at (3.6, 2.8) {$45^\circ$};
    \draw[quotientorange!80, thick] (3.3, 2.5) arc (0:-45:0.8);
    \node[quotientorange!90!black, font=\scriptsize] at (3.6, 2.2) {$-45^\circ$};
    
    \node[font=\scriptsize, align=center] at (2.5, -0.8) {Both $\ket{+}$ and $\ket{-}$ yield $50\%/50\%$ measurement outcomes,\\but they point in orthogonal geometric directions!};
  \end{scope}

  % Panel (b): Branching interference tree
  \begin{scope}[shift={(8.4, 0)}]
    \node[font=\small\bfseries] at (3.4, 6.1) {(b) Classical Probabilities vs. Quantum Amplitudes};
    
    % Classical Tree
    \node[draw=gray!60, rounded corners=2pt, font=\scriptsize\bfseries, fill=gray!10] at (1.2, 5.1) {Classical Bit $0$};
    \draw[->, thick, gray!70] (1.0, 4.8) -- (0.4, 3.9) node[midway, left, font=\tiny] {$p = \tfrac{1}{2}$};
    \draw[->, thick, gray!70] (1.4, 4.8) -- (2.0, 3.9) node[midway, right, font=\tiny] {$p = \tfrac{1}{2}$};
    
    \node[draw=gray!50, rounded corners=2pt, font=\tiny] (c0) at (0.4, 3.65) {State $0$};
    \node[draw=gray!50, rounded corners=2pt, font=\tiny] (c1) at (2.0, 3.65) {State $1$};
    
    \draw[->, thick, gray!70] (c0) -- (0.0, 2.6) node[midway, left, font=\tiny] {$\tfrac{1}{4}$};
    \draw[->, thick, gray!70] (c0) -- (0.8, 2.6) node[midway, right, font=\tiny] {$\tfrac{1}{4}$};
    \draw[->, thick, gray!70] (c1) -- (1.6, 2.6) node[midway, left, font=\tiny] {$\tfrac{1}{4}$};
    \draw[->, thick, gray!70] (c1) -- (2.4, 2.6) node[midway, right, font=\tiny] {$\tfrac{1}{4}$};
    
    \node[draw=gray!60, font=\tiny, fill=gray!5, align=left, inner sep=3pt] at (1.2, 1.8) {
      $\P(0) = \tfrac{1}{4}+\tfrac{1}{4} = \mathbf{50\%}$\\
      $\P(1) = \tfrac{1}{4}+\tfrac{1}{4} = \mathbf{50\%}$
    };
    \node[font=\tiny\itshape, text=gray!80!black, align=center] at (1.2, 0.9) {Probabilities only add ($p \ge 0$).\\Result stays random.};

    % Quantum Tree
    \node[draw=sourceblue!70, rounded corners=2pt, font=\scriptsize\bfseries, fill=sourceblue!10] at (5.4, 5.1) {Quantum Qubit $\ket{0}$};
    \draw[->, thick, sourceblue!70] (5.1, 4.8) -- (4.4, 3.9) node[midway, left, font=\tiny] {$\alpha = +\tfrac{1}{\sqrt{2}}$};
    \draw[->, thick, sourceblue!70] (5.7, 4.8) -- (6.4, 3.9) node[midway, right, font=\tiny] {$\alpha = +\tfrac{1}{\sqrt{2}}$};
    
    \node[draw=sourceblue!60, rounded corners=2pt, font=\tiny] (q0) at (4.4, 3.65) {Branch $\ket{0}$};
    \node[draw=sourceblue!60, rounded corners=2pt, font=\tiny] (q1) at (6.4, 3.65) {Branch $\ket{1}$};
    
    \draw[->, thick, targetgreen!80!black] (q0) -- (4.0, 2.6) node[midway, left, font=\tiny] {$+\tfrac{1}{2}$};
    \draw[->, thick, trapred!80!black] (q0) -- (4.8, 2.6) node[midway, right, font=\tiny] {$+\tfrac{1}{2}$};
    \draw[->, thick, targetgreen!80!black] (q1) -- (6.0, 2.6) node[midway, left, font=\tiny] {$+\tfrac{1}{2}$};
    \draw[->, thick, trapred!80!black] (q1) -- (6.8, 2.6) node[midway, right, font=\tiny] {$-\tfrac{1}{2}$};
    
    \node[draw=targetgreen!80!black, font=\tiny, fill=targetgreen!12, align=left, inner sep=3pt] at (5.4, 1.8) {
      $\ket{0}: +\tfrac{1}{2}+\tfrac{1}{2} = \mathbf{1} \implies \mathbf{100\%}$\\
      $\ket{1}: +\tfrac{1}{2}-\tfrac{1}{2} = \mathbf{0} \implies \mathbf{0\%}$
    };
    \node[font=\tiny\itshape, text=targetgreen!70!black, align=center] at (5.4, 0.9) {Amplitudes cancel ($+a - a = 0$).\\Interference restores certainty!};
  \end{scope}
\end{tikzpicture}
\caption{The foundational distinction of quantum computation. Panel (a) shows the single-qubit unit circle where $\ket{0}, \ket{1}$ are standard basis axes and $\ket{+}, \ket{-}$ are diagonal superpositions. Panel (b) compares two successive operations: two classical coin flips sum probabilities to maintain $50\%/50\%$ uncertainty, whereas two Hadamard gates cause destructive interference on $\ket{1}$ (canceling to $0$) and constructive interference on $\ket{0}$ (reinforcing to $100\%$).}
\label{ch42:fig-hadamard-interference}
\end{figure}

This is the central takeaway: the second Hadamard gate acts as an \textbf{interferometer}.
It takes a state that looked completely random under measurement and uses the hidden
phase to funnel all amplitude back into the original definite state.

\section{Multiple Qubits, Tensor Products, and Entanglement}
\label{ch42:multi-qubit}

A single qubit has two basis states, $\ket{0}$ and $\ket{1}$. What happens when we
assemble a register of multiple qubits?

\subsection{Two qubits: A 4D vector space}

Two classical bits can be in four possible configurations: $00$, $01$, $10$, or $11$.
A two-qubit quantum register has four corresponding standard basis vectors:
\[
  \ket{00} = \colvec{1\\0\\0\\0}, \quad
  \ket{01} = \colvec{0\\1\\0\\0}, \quad
  \ket{10} = \colvec{0\\0\\1\\0}, \quad
  \ket{11} = \colvec{0\\0\\0\\1}.
\]
An arbitrary two-qubit state is a unit vector in $\C^4$:
\[
  \ket{\psi} = \alpha_{00}\ket{00} + \alpha_{01}\ket{01} + \alpha_{10}\ket{10} + \alpha_{11}\ket{11},
\]
with normalization condition:
\[
  \abs{\alpha_{00}}^2 + \abs{\alpha_{01}}^2 + \abs{\alpha_{10}}^2 + \abs{\alpha_{11}}^2 = 1.
\]
When two independent qubits $\ket{\psi_1} = a\ket{0} + b\ket{1}$ and
$\ket{\psi_2} = c\ket{0} + d\ket{1}$ are placed next to each other, their joint state
is their \textbf{tensor product} (or Kronecker product):
\[
  \ket{\psi_1} \otimes \ket{\psi_2}
  = (a\ket{0} + b\ket{1}) \otimes (c\ket{0} + d\ket{1})
  = ac\ket{00} + ad\ket{01} + bc\ket{10} + bd\ket{11}.
\]
For $n$ qubits, the state space is $\C^{2^n}$. A 50-qubit processor requires $2^{50} \approx 10^{15}$
complex amplitudes to specify its complete state; a 300-qubit processor requires $2^{300}$
amplitudes, exceeding the number of particles in the universe.

\subsection{The CNOT Gate: Conditional interaction}

To perform logic between qubits, we need multi-qubit gates. The most important two-qubit
gate is the \textbf{Controlled-NOT (CNOT)} gate.

The CNOT gate acts on two qubits: a \textbf{control qubit} (the first qubit) and a
\textbf{target qubit} (the second qubit). If the control qubit is $0$, the target qubit
is left alone. If the control qubit is $1$, the target qubit is inverted ($X$ is applied):
\[
  \ket{00} \mapsto \ket{00}, \quad
  \ket{01} \mapsto \ket{01}, \quad
  \ket{10} \mapsto \ket{11}, \quad
  \ket{11} \mapsto \ket{10}.
\]
In matrix form acting on $\C^4$, CNOT is represented by:
\[
  \text{CNOT} = \begin{pmatrix}
    1 & 0 & 0 & 0 \\
    0 & 1 & 0 & 0 \\
    0 & 0 & 0 & 1 \\
    0 & 0 & 1 & 0
  \end{pmatrix}.
\]
Notice that CNOT is a permutation matrix: its rows and columns are orthonormal, so it is
unitary and reversible ($\text{CNOT}^\dagger = \text{CNOT}$ and $\text{CNOT}^2 = I$).

\subsection{Creating the Bell State: Entanglement demystified}

Now let us combine the Hadamard gate and the CNOT gate to produce the most famous state
in quantum physics: the \textbf{Bell State} (or EPR pair).

We start with two qubits initialized to $\ket{00}$.
\begin{enumerate}
  \item Apply $H$ to the first qubit:
  \[
    (H\ket{0}) \otimes \ket{0}
    = \left(\frac{1}{\sqrt{2}}\ket{0} + \frac{1}{\sqrt{2}}\ket{1}\right) \otimes \ket{0}
    = \frac{1}{\sqrt{2}}\ket{00} + \frac{1}{\sqrt{2}}\ket{10}.
  \]
  \item Apply CNOT with qubit 1 as control and qubit 2 as target:
  \begin{itemize}
    \item The branch $\ket{00}$ has control $0$, so target stays $0$: $\ket{00} \mapsto \ket{00}$.
    \item The branch $\ket{10}$ has control $1$, so target flips to $1$: $\ket{10} \mapsto \ket{11}$.
  \end{itemize}
\end{enumerate}
The resulting two-qubit state is:
\[
  \ket{\Phi^+} = \frac{1}{\sqrt{2}}\ket{00} + \frac{1}{\sqrt{2}}\ket{11}.
\]
This state is \textbf{entangled}. What does that word mean mathematically?

\begin{definition}[Entanglement]
A multi-qubit state $\ket{\psi} \in \C^{2^n}$ is \textbf{separable} (or a product state)
if it can be factored into individual single-qubit states:
$\ket{\psi} = \ket{\psi_1} \otimes \ket{\psi_2} \otimes \cdots \otimes \ket{\psi_n}$.
A state is \textbf{entangled} if no such factorization exists.
\end{definition}

Let us prove in one algebraic check that $\ket{\Phi^+}$ is entangled. Suppose for
contradiction that $\ket{\Phi^+}$ could be factored as:
\[
  \ket{\Phi^+} = (a\ket{0} + b\ket{1}) \otimes (c\ket{0} + d\ket{1})
  = ac\ket{00} + ad\ket{01} + bc\ket{10} + bd\ket{11}.
\]
Matching the coefficients of $\ket{\Phi^+} = \frac{1}{\sqrt{2}}\ket{00} + 0\ket{01} + 0\ket{10} + \frac{1}{\sqrt{2}}\ket{11}$ requires:
\[
  ac = \frac{1}{\sqrt{2}}, \qquad ad = 0, \qquad bc = 0, \qquad bd = \frac{1}{\sqrt{2}}.
\]
Look at $ad = 0$. If $a = 0$, then $ac = 0$, which contradicts $ac = 1/\sqrt{2}$.
Therefore $a \ne 0$, which forces $d = 0$. But if $d = 0$, then $bd = 0$, which contradicts
$bd = 1/\sqrt{2}$.

The contradiction is immediate. The state $\ket{\Phi^+}$ \textbf{cannot be factored}.
Neither individual qubit possesses a state of its own! The two-qubit system has a definite,
pure state, but the individual components do not.

\subsection{What entanglement does, and what it cannot do}

If Alice takes qubit 1 to Alpha Centauri and Bob keeps qubit 2 on Earth, what happens
when they measure?
\begin{itemize}
  \item If Alice measures qubit 1 in the standard basis, she sees $0$ with probability
  $(1/\sqrt{2})^2 = 50\%$, and $1$ with probability $(1/\sqrt{2})^2 = 50\%$.
  \item If Alice saw $0$, the state instantly collapses to $\ket{00}$. When Bob measures
  qubit 2, he is guaranteed to see $0$ with $100\%$ certainty!
  \item If Alice saw $1$, the state collapses to $\ket{11}$, and Bob is guaranteed to see $1$.
\end{itemize}
Their outcomes are $100\%$ correlated, no matter how far apart they are. As John Bell
proved in 1964, these correlations violate Bell's inequality, demonstrating that the
universe cannot be explained by local classical hidden variables.

However, \textbf{entanglement cannot transmit information faster than light}.
Notice that from Bob's local perspective, his qubit has probability $1/2$ of being $0$
and probability $1/2$ of being $1$. Alice cannot choose whether she gets $0$ or $1$; her
outcome is fundamentally random. Bob cannot tell whether Alice has measured her qubit or
not until Alice sends him a classical message (traveling at or below the speed of light)
telling him what she observed. Entanglement provides shared randomness with extraordinary
geometric correlations, but it never violates special relativity.

\section{Deutsch's Algorithm and the Phase-Kickback Trick}
\label{ch42:deutsch}

Now that we understand single-qubit interference and two-qubit gates, we can trace our
first complete quantum algorithm: David Deutsch's 1985 algorithm.

\subsection{The problem: Constant vs. balanced}

Suppose you are given a black-box function $f: \{0, 1\} \to \{0, 1\}$ that takes a
single bit as input and outputs a single bit. There are four such functions:
\begin{enumerate}
  \item $f(0) = 0$ and $f(1) = 0$ (Constant)
  \item $f(0) = 1$ and $f(1) = 1$ (Constant)
  \item $f(0) = 0$ and $f(1) = 1$ (Balanced)
  \item $f(0) = 1$ and $f(1) = 0$ (Balanced)
\end{enumerate}
A function is \textbf{constant} if $f(0) = f(1)$, and \textbf{balanced} if $f(0) \ne f(1)$.

\textbf{The classical query complexity:}
How many times must a classical computer evaluate $f$ to determine with certainty
whether $f$ is constant or balanced?
The answer is strictly \textbf{two}. If you evaluate $f(0)$ and obtain $0$, you know
nothing about whether $f$ is constant or balanced; $f(1)$ could be $0$ (constant) or $1$
(balanced). You must evaluate both $f(0)$ and $f(1)$.

Can a quantum computer solve this problem in \textbf{one single query}?

\subsection{The quantum oracle and Phase Kickback}

In quantum mechanics, all gate operations must be reversible. If $f$ is a constant
function, the mapping $\ket{x} \mapsto \ket{f(x)}$ erases the input $x$ and is not
unitary. To make evaluation reversible, we encode $f$ into a two-qubit unitary oracle
$U_f$:
\[
  U_f \ket{x}\ket{y} = \ket{x}\ket{y \oplus f(x)},
\]
where $\oplus$ denotes addition modulo $2$ (bitwise XOR). The first qubit $\ket{x}$ is the
input, and the second qubit $\ket{y}$ is the target scratchpad.

Now we introduce the fundamental building block of quantum algorithms:
\textbf{Phase Kickback}.

What happens if we prepare the target qubit in the state $\ket{-} = \frac{1}{\sqrt{2}}(\ket{0} - \ket{1})$?
Let us calculate $U_f \ket{x}\ket{-}$ explicitly:
\begin{align*}
  U_f \ket{x}\ket{-} &= U_f \ket{x}\left(\frac{\ket{0} - \ket{1}}{\sqrt{2}}\right) \\
  &= \frac{1}{\sqrt{2}}\Big(\ket{x}\ket{0 \oplus f(x)} - \ket{x}\ket{1 \oplus f(x)}\Big).
\end{align*}
Let us check the two possible values of $f(x)$:
\begin{itemize}
  \item If $f(x) = 0$:
  \[
    \frac{1}{\sqrt{2}}\Big(\ket{x}\ket{0} - \ket{x}\ket{1}\Big)
    = \ket{x}\left(\frac{\ket{0} - \ket{1}}{\sqrt{2}}\right) = +\ket{x}\ket{-}.
  \]
  \item If $f(x) = 1$:
  \[
    \frac{1}{\sqrt{2}}\Big(\ket{x}\ket{1} - \ket{x}\ket{0}\Big)
    = -\ket{x}\left(\frac{\ket{0} - \ket{1}}{\sqrt{2}}\right) = -\ket{x}\ket{-}.
  \]
\end{itemize}
In both cases, we can write the output compactly as:
\[
  \boxed{U_f \ket{x}\ket{-} = (-1)^{f(x)}\ket{x}\ket{-}.}
\]
The target qubit $\ket{-}$ emerged completely unchanged! But its presence caused the value
$f(x)$ to be \textbf{kicked back} into the phase of the input qubit as a factor of $(-1)^{f(x)}$.

\subsection{Tracing Deutsch's algorithm step-by-step}

Here is the complete quantum circuit:
\begin{enumerate}
  \item Initialize the register to $\ket{\psi_0} = \ket{0}\ket{1}$.
  \item Apply Hadamard to both qubits:
  \[
    \ket{\psi_1} = (H\ket{0}) \otimes (H\ket{1})
    = \left(\frac{\ket{0} + \ket{1}}{\sqrt{2}}\right) \otimes \ket{-}.
  \]
  \item Query the oracle $U_f$ once:
  \begin{align*}
    \ket{\psi_2} &= U_f \left(\frac{\ket{0} + \ket{1}}{\sqrt{2}} \otimes \ket{-}\right) \\
    &= \frac{1}{\sqrt{2}} U_f\ket{0}\ket{-} + \frac{1}{\sqrt{2}} U_f\ket{1}\ket{-} \\
    &= \frac{1}{\sqrt{2}}\Big((-1)^{f(0)}\ket{0} + (-1)^{f(1)}\ket{1}\Big) \otimes \ket{-}.
  \end{align*}
  \item Factor out $(-1)^{f(0)}$:
  \[
    \ket{\psi_2} = (-1)^{f(0)} \left(\frac{\ket{0} + (-1)^{f(0)\oplus f(1)}\ket{1}}{\sqrt{2}}\right) \otimes \ket{-}.
  \]
  Now look at the term $(-1)^{f(0)\oplus f(1)}$:
  \begin{itemize}
    \item \textbf{If $f$ is constant:} $f(0) = f(1)$, so $f(0)\oplus f(1) = 0$. The first qubit is:
    \[
      \pm \left(\frac{\ket{0} + \ket{1}}{\sqrt{2}}\right) = \pm \ket{+}.
    \]
    \item \textbf{If $f$ is balanced:} $f(0) \ne f(1)$, so $f(0)\oplus f(1) = 1$. The first qubit is:
    \[
      \pm \left(\frac{\ket{0} - \ket{1}}{\sqrt{2}}\right) = \pm \ket{-}.
    \]
  \end{itemize}
  \item Apply Hadamard to the first qubit:
  Because $H\ket{+} = \ket{0}$ and $H\ket{-} = \ket{1}$, the first qubit becomes:
  \[
    \ket{\psi_3} = \begin{cases}
      \pm\ket{0} & \text{if } f \text{ is constant}, \\
      \pm\ket{1} & \text{if } f \text{ is balanced}.
    \end{cases}
  \]
  \item Measure the first qubit:
  If the outcome is $0$, $f$ is constant. If the outcome is $1$, $f$ is balanced.
\end{enumerate}

The algorithm determines whether $f$ is constant or balanced with $100\%$ certainty
in \textbf{a single query}.

How was this possible? Deutsch's algorithm did not learn $f(0)$, and it did not learn
$f(1)$. If you asked the computer what $f(0)$ was, it would have no idea! Instead, it
used quantum interference to evaluate the \textbf{global relational property}
$f(0)\oplus f(1)$ directly, canceling the individual values while preserving their
relative phase.

\section{Grover's Search as 2D Plane Geometry}
\label{ch42:grover}

Deutsch's algorithm demonstrated a modest advantage ($1$ query vs. $2$). In 1996, Lov
Grover discovered that quantum interference could accelerate unstructured database
search across $N = 2^n$ items, finding a unique marked item in $\Theta(\sqrt{N})$
steps where classical search requires $\Theta(N)$ steps.

Most textbooks explain Grover's algorithm using pages of algebraic summation. Here
we give the 3Blue1Brown geometric explanation: Grover's search is simply a
\textbf{rotation in a two-dimensional plane}.

\subsection{The 2D Grover plane}

Suppose we have an unsorted database of $N = 2^n$ elements, represented by $n$-bit
strings $x \in \{0, 1\}^n$. Exactly one marked item $x^\star$ satisfies $f(x^\star) = 1$,
while $f(x) = 0$ for all $x \ne x^\star$.

Instead of attempting to visualize a $2^n$-dimensional space, notice that the entire
algorithm takes place in a flat two-dimensional subspace spanned by just two vectors:
\begin{enumerate}
  \item The target state: $\ket{x^\star}$.
  \item The uniform superposition of all $N-1$ non-target states:
  \[
    \ket{w} = \frac{1}{\sqrt{N-1}}\sum_{x \ne x^\star}\ket{x}.
  \]
\end{enumerate}
Notice that $\braket{w}{x^\star} = 0$: the two vectors are strictly orthogonal! We can
treat $\ket{w}$ as our horizontal $x$-axis and $\ket{x^\star}$ as our vertical $y$-axis.

We initialize our algorithm in the uniform superposition over all $N$ items:
\[
  \ket{s} = \frac{1}{\sqrt{N}}\sum_{x=0}^{N-1}\ket{x}
  = \sqrt{\frac{N-1}{N}}\ket{w} + \frac{1}{\sqrt{N}}\ket{x^\star}.
\]
Define the angle $\theta$ such that $\sin\theta = \frac{1}{\sqrt{N}}$ and
$\cos\theta = \sqrt{\frac{N-1}{N}}$. Then:
\[
  \ket{s} = \cos\theta\ket{w} + \sin\theta\ket{x^\star}.
\]
When $N$ is large (say $N = 10^6$), $\sin\theta \approx 1/\sqrt{N} = 1/1000$, which
means $\theta$ is a tiny angle. The initial vector $\ket{s}$ lies almost flat against the
horizontal $\ket{w}$-axis. Our goal is to rotate this arrow until it points vertically
along $\ket{x^\star}$.

\subsection{Two reflections make a rotation}

A famous theorem of high-school plane geometry states:
\begin{theorem}[Composition of Reflections]
Reflecting a vector across a line $L_1$, and then reflecting the result across a line
$L_2$ that intersects $L_1$ at an angle $\theta$, rotates the vector by an angle of
$2\theta$ in the direction from $L_1$ toward $L_2$.
\end{theorem}

Grover's algorithm implements exactly this theorem! Each Grover iteration consists of
two operations:
\begin{enumerate}
  \item \textbf{Operation 1: The Oracle Reflection ($U_f$).}
  The oracle flips the sign of the target state while leaving all other states alone:
  $U_f = I - 2\ket{x^\star}\bra{x^\star}$.
  In our 2D plane:
  \[
    U_f\ket{w} = \ket{w}, \qquad U_f\ket{x^\star} = -\ket{x^\star}.
  \]
  Geometrically, this flips the vertical component of the vector while leaving the
  horizontal component unchanged. This is a \textbf{reflection across the horizontal
  axis $\ket{w}$}!
  
  When applied to our state $\ket{s}$ (at angle $+\theta$), the oracle reflects it to
  angle $-\theta$.

  \item \textbf{Operation 2: The Diffusion Reflection ($U_s$).}
  The diffusion operator is defined by:
  \[
    U_s = 2\ket{s}\bra{s} - I.
  \]
  If you take a vector parallel to $\ket{s}$, $U_s\ket{s} = 2\ket{s}(1) - \ket{s} = +\ket{s}$.
  If you take a vector orthogonal to $\ket{s}$, $U_s\ket{s_\perp} = 2\ket{s}(0) - \ket{s_\perp} = -\ket{s_\perp}$.
  This is a \textbf{reflection across the line along $\ket{s}$}!
\end{enumerate}

Now put the two reflections together:
\begin{itemize}
  \item The state was at angle $+\theta$.
  \item Oracle reflection across line $L_1$ (horizontal axis at angle $0$) sent it to angle $-\theta$.
  \item The angular distance from the new state ($-\theta$) to line $L_2$ (vector $\ket{s}$ at angle $+\theta$) is $2\theta$.
  \item Diffusion reflection across line $L_2$ reflects it across $\ket{s}$, placing it at angle:
  \[
    \theta + 2\theta = 3\theta!
  \]
\end{itemize}
In one single Grover iteration $G = U_s U_f$, the state arrow has rotated by an angle of
\textbf{exactly $2\theta$} toward the vertical target axis $\ket{x^\star}$!

\subsection{How many steps to find the needle?}

Every iteration advances the angle by $2\theta$. We started at angle $\theta$. To point
straight up at $\ket{x^\star}$, we need the angle to reach $\pi/2$ ($90^\circ$).
The remaining angular distance is:
\[
  \frac{\pi}{2} - \theta \approx \frac{\pi}{2}.
\]
Because each iteration adds $2\theta \approx \frac{2}{\sqrt{N}}$ radians, the total number
of iterations $R$ required is:
\[
  R \approx \frac{\pi/2}{2\theta} \approx \frac{\pi/2}{2/\sqrt{N}} = \frac{\pi}{4}\sqrt{N}.
\]

\begin{figure}[H]
\centering
\begin{tikzpicture}
  % Panel (a): 2D Grover plane reflections
  \begin{scope}[shift={(0,0)}]
    \node[font=\small\bfseries] at (2.7, 5.5) {(a) Grover Step as Two Plane Reflections};
    
    % Axes
    \draw[->, gray!70, thick] (-0.3, 0) -- (5.2, 0) node[right, font=\footnotesize] {$\ket{w}$ (non-targets)};
    \draw[->, gray!70, thick] (0, -1.8) -- (0, 4.6) node[above, font=\footnotesize] {$\ket{x^\star}$ (target)};
    
    % Vectors
    \draw[->, very thick, sourceblue!90!black] (0,0) -- (4.25, 1.14) node[above right, font=\footnotesize] {$\ket{s}$ (start)};
    \draw[->, very thick, trapred!80!black] (0,0) -- (4.25, -1.14) node[below right, font=\footnotesize] {$U_f\ket{s}$ (after oracle)};
    \draw[->, very thick, targetgreen!85!black] (0,0) -- (3.11, 3.11) node[above right, font=\footnotesize] {$U_s U_f\ket{s}$ (after diffusion)};
    
    % Reflection arrows
    \draw[->, trapred, thick, dashed] (3.8, 0.9) -- (3.8, -0.9) node[pos=0.72, left=2pt, font=\tiny] {1. Reflect across $\ket{w}$};
    \draw[->, targetgreen!80!black, thick, dashed] (4.0, -1.0) arc (-15:45:4.1) node[pos=0.72, right=2pt, font=\tiny] {2. Reflect across $\ket{s}$};
    
    % Angle arcs
    \draw[sourceblue!80, thick] (1.6, 0) arc (0:15:1.6);
    \node[sourceblue!90!black, font=\tiny] at (1.9, 0.23) {$\theta$};
    \draw[trapred!80, thick] (1.6, 0) arc (0:-15:1.6);
    \node[trapred!90!black, font=\tiny] at (1.9, -0.23) {$\theta$};
    \draw[targetgreen!80!black, thick] (1.3, 0.35) arc (15:45:1.3);
    \node[targetgreen!90!black, font=\tiny] at (1.7, 0.92) {$2\theta$};
    
    \node[font=\scriptsize, align=center] at (2.7, -2.4) {Each step rotates the state by $2\theta \approx 2/\sqrt{N}$.\\After $\approx \tfrac{\pi}{4}\sqrt{N}$ steps, the arrow points at $\ket{x^\star}$.};
  \end{scope}

  % Panel (b): Inversion about the mean bar chart
  \begin{scope}[shift={(8.2, 0)}]
    \node[font=\small\bfseries] at (3.4, 5.5) {(b) Amplitude Inversion About the Mean};
    
    % Stage 1: Uniform
    \node[anchor=west, font=\scriptsize\bfseries] at (0.0, 4.4) {1. Initial Superposition:};
    \draw[->, gray!60] (0, 3.4) -- (6.6, 3.4);
    \foreach \i in {0.6, 1.4, 2.2, 3.0, 3.8, 5.4} {
      \fill[sourceblue!70] (\i, 3.4) rectangle ++(0.35, 0.55);
    }
    \node[targetgreen!80!black, font=\tiny] at (5.57, 4.18) {$x^\star$};
    \draw[dashed, quotientorange, thick] (0.2, 3.95) -- (6.3, 3.95) node[right, font=\tiny] {$\mu = \tfrac{1}{\sqrt{N}}$};

    % Stage 2: After Oracle
    \node[anchor=west, font=\scriptsize\bfseries] at (0.0, 2.55) {2. After Oracle $U_f$ (Sign Flip):};
    \draw[->, gray!60] (0, 1.65) -- (6.6, 1.65);
    \foreach \i in {0.6, 1.4, 2.2, 3.0, 3.8} {
      \fill[sourceblue!70] (\i, 1.65) rectangle ++(0.35, 0.55);
    }
    \fill[trapred!80] (5.4, 1.65) rectangle ++(0.35, -0.55);
    \node[trapred!80!black, font=\tiny, right=2pt] at (5.75, 1.35) {$-1/\sqrt{N}$};
    \draw[dashed, quotientorange, thick] (0.2, 2.05) -- (6.3, 2.05) node[right, font=\tiny] {Mean $\mu$};

    % Stage 3: After Diffusion
    \node[anchor=west, font=\scriptsize\bfseries] at (0.0, 0.45) {3. After Diffusion $2\ket{s}\bra{s}-I$ (Invert about mean):};
    \draw[->, gray!60] (0, -0.95) -- (6.6, -0.95);
    \foreach \i in {0.6, 1.4, 2.2, 3.0, 3.8} {
      \fill[sourceblue!40] (\i, -0.95) rectangle ++(0.35, 0.2);
    }
    \fill[targetgreen!85!black] (5.4, -0.95) rectangle ++(0.35, 1.15);
    \node[targetgreen!85!black, font=\tiny\bfseries, right=2pt] at (5.75, -0.1) {Boosted!};
    \draw[dashed, quotientorange, thick] (0.2, -0.7) -- (6.3, -0.7) node[right, font=\tiny] {$\mu$};
    
    \node[font=\scriptsize, align=center] at (3.4, -2.4) {Diffusion maps $\alpha_x \mapsto 2\mu - \alpha_x$.\\The negative target amplitude reflects high above the mean!};
  \end{scope}
\end{tikzpicture}
\caption{Grover's algorithm as geometric rotation and amplitude amplification. In panel (a), the oracle reflects the state across $\ket{w}$, and the diffusion operator reflects across $\ket{s}$, combining to produce an exact rotation of $2\theta \approx 2/\sqrt{N}$ per iteration. In panel (b), the algebraic diffusion operation $2\ket{s}\bra{s}-I$ inverts every amplitude about the mean $\mu$, amplifying the negative target amplitude into a tall positive peak.}
\label{ch42:fig-grover-geometry}
\end{figure}

\begin{trapbox}
Unlike classical algorithms that improve monotonically the longer they run, Grover's
algorithm is a \textbf{rotation}. If you continue applying iterations past
$R \approx \frac{\pi}{4}\sqrt{N}$, the arrow rotates right past $\ket{x^\star}$ and
points back toward the horizontal axis! The probability of success drops toward zero.
You must stop the machine and measure at precisely the right number of steps.
\end{trapbox}

\section{Shor's Algorithm, the BQP Landscape, and Fault Tolerance}
\label{ch42:shor-bqp}

Grover's search offers a quadratic speedup for unstructured problems. Can quantum
computers achieve an \textbf{exponential speedup}?

\subsection{Shor's algorithm and the Quantum Fourier Transform}

In 1994, Peter Shor answered with integer factorization. Given a composite number
$N = pq$, Shor's algorithm finds its prime factors in $O((\log N)^3)$ operations,
compared to the best known classical algorithm (the General Number Field Sieve), which
runs in sub-exponential time $\exp(O((\log N)^{1/3}(\log\log N)^{2/3}))$.

Shor's breakthrough connects directly to the discrete Fourier transform and the FFT
developed in Chapter 29. The algorithm proceeds in two stages:
\begin{enumerate}
  \item \textbf{Classical Reduction to Period Finding:}
  Pick an integer $a$ coprime to $N$ ($1 < a < N$). The sequence of modular powers
  $f(x) = a^x \bmod N$ is periodic: there exists a smallest positive integer $r$ (the
  order of $a$) such that $a^r \equiv 1 \pmod N$.
  If $r$ is even and $a^{r/2} \not\equiv -1 \pmod N$, then
  $(a^{r/2}-1)(a^{r/2}+1) = a^r - 1 \equiv 0 \pmod N$.
  By Euclid's algorithm, $\gcd(a^{r/2}-1, N)$ reveals a non-trivial prime factor of $N$!
  Classically, finding $r$ requires testing exponentially many values of $x$.

  \item \textbf{Quantum Period Extraction via the QFT:}
  A quantum register is prepared in an equal superposition over inputs $x \in \{0, \dots, Q-1\}$,
  and the modular exponentiation $a^x \bmod N$ is evaluated in a second register.
  Measuring the second register collapses the first register into a periodic spike train:
  superposition over all inputs $x_0, x_0 + r, x_0 + 2r, \dots$ spaced by the unknown period $r$.
  
  Applying the \textbf{Quantum Fourier Transform (QFT)} turns this time-domain periodic
  train into frequency-domain interference peaks concentrated at integer multiples of
  $Q/r$. Measuring the output reveals a frequency peak, from which the period $r$ is
  extracted in polynomial time using continued fractions.
\end{enumerate}

Shor's algorithm succeeds exponentially because integer factorization possesses rich
underlying algebraic symmetry (the periodic structure of abelian modular groups).

\subsection{Does a quantum computer solve NP-complete problems?}

A frequent question from computer scientists is:
\begin{quote}
\emph{``Can a quantum computer solve NP-complete problems (like 3-SAT, Traveling Salesperson,
or Clique) in polynomial time?''}
\end{quote}

The consensus among complexity theorists is: \textbf{almost certainly not}.

NP-complete problems do not possess the global periodic group structure that Shor's QFT
exploits. An instance of 3-SAT with $n$ variables has $2^n$ possible truth assignments,
with no algebraic periodicity linking satisfying solutions. Against an unstructured
combinatorial landscape, the most effective generic quantum strategy known is Grover's
search.

Grover's algorithm accelerates searching the $2^n$ assignments from $O(2^n)$ down to
$O(2^{n/2})$. While taking the square root of the search space is a remarkable mathematical
achievement, $2^{n/2}$ is still an \textbf{exponential runtime}. For a 3-SAT instance
with $n = 200$ boolean variables:
\[
  2^{n/2} = 2^{100} \approx 1.27 \times 10^{30} \text{ operations.}
\]
A computer running a billion iterations every microsecond would still require tens of
thousands of years to complete the search. Quantum computers do not magically eliminate
combinatorial explosion.

\subsection{The Complexity Landscape: Situating BQP}

To formalize quantum computational power, Bernstein and Vazirani defined the complexity
class $\mathrm{BQP}$:

\begin{definition}[BQP]
$\mathrm{BQP}$ (Bounded-Error Quantum Polynomial Time) is the class of decision problems
solvable by a uniform family of polynomial-size quantum circuits with an error
probability at most $1/3$ on all inputs.
\end{definition}

How does $\mathrm{BQP}$ compare to the classical complexity classes studied in Chapter 33?
\begin{itemize}
  \item $\mathrm{P} \subseteq \mathrm{BPP} \subseteq \mathrm{BQP} \subseteq \mathrm{PSPACE}$.
  Every problem efficiently solvable by classical computers (deterministic or randomized)
  is efficiently solvable by a quantum computer.
  \item \textbf{Factoring and Discrete Logarithms} are in $\mathrm{BQP}$. Classically,
  they reside in $\mathrm{NP} \cap \text{co-NP}$, but are not believed to be in $\mathrm{P}$.
  \item \textbf{Quantum Simulation} (simulating the time evolution of local quantum
  Hamiltonians) is in $\mathrm{BQP}$, and is $\mathrm{BQP}$-complete.
  \item $\mathrm{BQP}$ is \textbf{not} believed to contain $\mathrm{NP}$-complete problems.
  Indeed, there are oracle worlds where $\mathrm{BQP}$ and $\mathrm{NP}$ are incomparable:
  $\mathrm{BQP}$ contains problems outside $\mathrm{NP}$, and $\mathrm{NP}$ contains
  problems outside $\mathrm{BQP}$.
\end{itemize}

\begin{figure}[H]
\centering
\begin{tikzpicture}
  % Panel (a): Complexity Venn Diagram
  \begin{scope}[shift={(0,0)}]
    \node[font=\small\bfseries] at (3.0, 5.2) {(a) The Complexity Landscape};
    
    % PSPACE boundary
    \draw[rounded corners=8pt, thick, gray!80, fill=gray!5] (-0.3, -0.4) rectangle (6.3, 4.8);
    \node[gray!80!black, font=\small\bfseries] at (0.6, 4.4) {$\mathrm{PSPACE}$};
    
    % NP Ellipse
    \draw[rotate around={-15:(3.2, 2.2)}, thick, quotientorange!80!black, fill=quotientorange!12] (3.2, 2.2) ellipse (2.7 and 1.8);
    \node[quotientorange!90!black, font=\small\bfseries] at (1.1, 3.4) {$\mathrm{NP}$};
    
    % BQP Ellipse
    \draw[rotate around={20:(3.0, 1.8)}, thick, sourceblue!80!black, fill=sourceblue!15, fill opacity=0.8] (3.0, 1.8) ellipse (2.6 and 1.5);
    \node[sourceblue!90!black, font=\small\bfseries] at (5.2, 2.8) {$\mathrm{BQP}$};
    
    % P Circle inside NP and BQP
    \draw[thick, targetgreen!80!black, fill=targetgreen!20] (2.2, 1.3) circle (0.85);
    \node[targetgreen!90!black, font=\small\bfseries] at (2.2, 1.3) {$\mathrm{P}$};
    
    % Problem markers
    \node[font=\tiny\bfseries, align=center] at (2.2, 0.7) {Sorting, Shortest Paths};
    \node[font=\tiny\bfseries, text=trapred!90!black, align=center] at (1.8, 3.1) {$\mathrm{NP}$-Complete\\(3-SAT, TSP, Clique)};
    \node[font=\tiny\bfseries, text=structurepurple!90!black, align=center] at (3.5, 2.2) {Factoring, Discrete Log\\(in $\mathrm{BQP} \cap \mathrm{NP}$, not in $\mathrm{P}$)};
    \node[font=\tiny\bfseries, text=sourceblue!90!black, align=center] at (4.2, 1.0) {Quantum Simulation\\($\mathrm{BQP}$-Complete)};
  \end{scope}

  % Panel (b): Surface Code Parity Check Grid
  \begin{scope}[shift={(7.5, 0)}]
    \node[font=\small\bfseries] at (3.0, 5.2) {(b) Surface Code Syndrome Measurements};
    
    % Grid lines
    \draw[gray!40, thick] (0.8, 0.4) grid[step=1.6] (4.0, 3.6);
    
    % Plaquettes
    \fill[targetgreen!25] (0.8, 2.0) rectangle ++(1.6, 1.6);
    \node[targetgreen!80!black, font=\scriptsize\bfseries] at (1.6, 2.8) {$Z$-Check ($Z^{\otimes 4}$)};
    \node[targetgreen!70!black, font=\tiny] at (1.6, 2.5) {Detects Bit-Flips};
    
    \fill[quotientorange!25] (2.4, 0.4) rectangle ++(1.6, 1.6);
    \node[quotientorange!80!black, font=\scriptsize\bfseries] at (3.2, 1.2) {$X$-Check ($X^{\otimes 4}$)};
    \node[quotientorange!70!black, font=\tiny] at (3.2, 0.9) {Detects Phase-Flips};

    \fill[quotientorange!25] (2.4, 2.0) rectangle ++(1.6, 1.6);
    \node[quotientorange!80!black, font=\scriptsize\bfseries] at (3.2, 2.8) {$X$-Check};
    
    \fill[targetgreen!25] (0.8, 0.4) rectangle ++(1.6, 1.6);
    \node[targetgreen!80!black, font=\scriptsize\bfseries] at (1.6, 1.2) {$Z$-Check};
    
    % Data qubits (vertices)
    \foreach \x in {0.8, 2.4, 4.0} {
      \foreach \y in {0.4, 2.0, 3.6} {
        \filldraw[fill=white, draw=gray!80!black, thick] (\x, \y) circle (4pt);
      }
    }
    \node[font=\tiny] at (0.8, 3.9) {$q_1$};
    \node[font=\tiny] at (2.4, 3.9) {$q_2$};
    \node[font=\tiny] at (4.0, 3.9) {$q_3$};
    \node[font=\tiny] at (0.8, 0.1) {$q_7$};
    \node[font=\tiny] at (2.4, 0.1) {$q_8$};
    \node[font=\tiny] at (4.0, 0.1) {$q_9$};
    
    \node[font=\footnotesize, align=center] at (2.4, -0.7) {Measuring $Z_1 Z_2 Z_4 Z_5$ checks mutual parity ($+1$ or $-1$)\\without measuring individual qubits, preserving superposition!};
  \end{scope}
\end{tikzpicture}
\caption{The theoretical and physical landscape of quantum advantage. Panel (a) illustrates the complexity hierarchy: $\mathrm{BQP}$ encompasses $\mathrm{P}$ and contains problems like Factoring and Quantum Simulation, while $\mathrm{NP}$-complete problems remain outside. Panel (b) depicts Kitaev's topological surface code, where physical data qubits reside on vertices and alternating plaquettes perform non-destructive parity checks ($Z$-syndromes for bit-flips, $X$-syndromes for phase-flips) to protect quantum information indefinitely.}
\label{ch42:fig-bqp-landscape}
\end{figure}

\subsection{Quantum Error Correction and the No-Cloning Theorem}

In classical computing, error correction is conceptually simple: if noise corrupts a
bit, you can back up the data by making redundant copies ($0 \mapsto 000$). If one bit
flips, majority voting restores the truth.

In quantum computing, this simple duplication is physically impossible:

\begin{theorem}[The No-Cloning Theorem, Wootters \& Zurek, 1982]
There is no unitary operator $U$ that can create an identical duplicate of an
unknown arbitrary quantum state:
\[
  U(\ket{\psi}\ket{0}) = \ket{\psi}\ket{\psi} \quad \text{for all } \ket{\psi}.
\]
\end{theorem}

\begin{proof}
Suppose such a unitary operator $U$ existed. Then $U$ must clone the basis states:
\[
  U(\ket{0}\ket{0}) = \ket{0}\ket{0}, \qquad U(\ket{1}\ket{0}) = \ket{1}\ket{1}.
\]
Now consider the superposition state $\ket{+} = \frac{1}{\sqrt{2}}(\ket{0} + \ket{1})$.
Because $U$ is a linear operator, applying $U$ to $\ket{+}\ket{0}$ must act linearly:
\[
  U\left(\frac{\ket{0} + \ket{1}}{\sqrt{2}} \otimes \ket{0}\right)
  = \frac{1}{\sqrt{2}}U(\ket{0}\ket{0}) + \frac{1}{\sqrt{2}}U(\ket{1}\ket{0})
  = \frac{1}{\sqrt{2}}\ket{00} + \frac{1}{\sqrt{2}}\ket{11}.
\]
However, if $U$ successfully cloned the state $\ket{+}$, the output would be:
\[
  \ket{+}\ket{+} = \left(\frac{\ket{0} + \ket{1}}{\sqrt{2}}\right) \otimes \left(\frac{\ket{0} + \ket{1}}{\sqrt{2}}\right)
  = \frac{1}{2}\ket{00} + \frac{1}{2}\ket{01} + \frac{1}{2}\ket{10} + \frac{1}{2}\ket{11}.
\]
These two states are manifestly unequal. A linear operator cannot clone non-orthogonal states.
\end{proof}

If you cannot clone a qubit, and looking at a qubit collapses its superposition, how
can quantum errors ever be corrected?

The insight of Peter Shor and Alexei Kitaev was to encode one \textbf{logical qubit}
into an entangled state across multiple physical qubits. Consider the 3-qubit
bit-flip code:
\[
  \ket{0_L} = \ket{000}, \qquad \ket{1_L} = \ket{111}.
\]
An arbitrary logical state is $\alpha\ket{000} + \beta\ket{111}$.

Instead of measuring the individual qubits (which would reveal whether the state is $0$
or $1$ and destroy $\alpha$ and $\beta$), we measure \textbf{parity-check operators}
using helper (ancilla) qubits:
\begin{itemize}
  \item Measure the parity $Z_1 Z_2$: returns $+1$ if qubits 1 and 2 agree, and $-1$ if they differ.
  \item Measure the parity $Z_2 Z_3$: returns $+1$ if qubits 2 and 3 agree, and $-1$ if they differ.
\end{itemize}
If an environmental disturbance flips qubit 1 ($X_1$), the parities read
$(Z_1 Z_2 = -1, Z_2 Z_3 = +1)$. The syndrome identifies precisely which qubit flipped,
allowing a corrective $X_1$ gate to restore the state, \textbf{without the measurement
ever learning whether the data was $\ket{000}$ or $\ket{111}$}!

In modern architectures, this principle is scaled to two-dimensional grids of qubits called
\textbf{surface codes} (Figure~\ref{ch42:fig-bqp-landscape}b). Data qubits sit on the vertices
of a lattice, and interleaved measurement qubits constantly check local parities to detect
bit-flips and phase-flips. By scaling the distance of the code, error rates can be made
arbitrarily small, bridging the gap between fragile quantum physics and fault-tolerant
computer science.

\begin{aestheticbox}
Quantum computing does not replace classical algorithms; it expands the physical grammar
of what an algorithm can be. By replacing positive probabilities with complex amplitudes,
it transforms computation from a sequence of irreversible decisions into a geometric
wave of interfering possibilities. Where classical algorithms must choose a branch,
quantum algorithms choreograph the paths so that false branches destroy themselves,
leaving the true answer to step forward into the light of measurement.
\end{aestheticbox}

\section{Exercises}

\begin{exercise}\label{ch42:ex-hadamard-trace}
\textbf{(Hadamard interference by hand.)}
Suppose a single qubit begins in the state $\ket{1} = \colvec{0\\1}$.
\begin{enumerate}
  \item Compute the state vector after applying the Hadamard gate: $\ket{\psi_1} = H\ket{1}$.
  \item What is the probability of observing $0$ if $\ket{\psi_1}$ is measured? What is the probability of observing $1$?
  \item Compute the state vector after applying a second Hadamard gate: $\ket{\psi_2} = H\ket{\psi_1}$. Write down each of the four paths from the intermediate basis states to the final basis states, and verify that the paths to $\ket{0}$ cancel to exact zero while the paths to $\ket{1}$ reinforce.
\end{enumerate}
\end{exercise}

\begin{exercise}\label{ch42:ex-bell-factorization}
\textbf{(Entanglement of the Bell states.)}
The four Bell states form an orthonormal basis for two qubits:
\[
  \ket{\Phi^\pm} = \frac{\ket{00} \pm \ket{11}}{\sqrt{2}}, \qquad
  \ket{\Psi^\pm} = \frac{\ket{01} \pm \ket{10}}{\sqrt{2}}.
\]
\begin{enumerate}
  \item Prove that $\ket{\Psi^+}$ cannot be factored as a product state
  $(\alpha_1\ket{0} + \beta_1\ket{1}) \otimes (\alpha_2\ket{0} + \beta_2\ket{1})$.
  \item Verify that the four Bell states are mutually orthogonal: show that
  $\braket{\Phi^+}{\Phi^-} = 0$ and $\braket{\Phi^+}{\Psi^+} = 0$.
\end{enumerate}
\end{exercise}

\begin{exercise}\label{ch42:ex-deutsch-trace}
\textbf{(Deutsch's algorithm with an explicit balanced function.)}
Let $f: \{0, 1\} \to \{0, 1\}$ be the balanced function $f(x) = x \oplus 1$ (the classical NOT function, so $f(0) = 1$ and $f(1) = 0$).
Trace Deutsch's algorithm step-by-step for this specific function:
\begin{enumerate}
  \item Write down the state after initializing $\ket{0}\ket{1}$ and applying $H \otimes H$.
  \item Calculate the state after applying the oracle $U_f$. Explicitly show how the phase kickback $(-1)^{f(x)}$ produces a minus sign on the $\ket{0}$ branch and a plus sign on the $\ket{1}$ branch of the first qubit.
  \item Apply $H$ to the first qubit and show that the final measured state is $\ket{1}$ with probability $1$, confirming that $f$ is balanced.
\end{enumerate}
\end{exercise}

\begin{exercise}\label{ch42:ex-grover-n4}
\textbf{(Grover's search on four items.)}
Consider a database of $N = 4$ items ($n = 2$ qubits), with marked target $x^\star = 11$.
\begin{enumerate}
  \item What is the initial state $\ket{s}$ in terms of the standard basis vectors $\ket{00}, \ket{01}, \ket{10}, \ket{11}$?
  \item Calculate the angle $\theta$ such that $\sin\theta = 1/\sqrt{N} = 1/2$. What is $\theta$ in degrees?
  \item Recall that one Grover iteration rotates the vector by $2\theta$. Compute the angle of the state after exactly $R = 1$ Grover iteration.
  \item Show that for $N = 4$, exactly \textbf{one single Grover iteration} rotates the state to an angle of $90^\circ$ ($\pi/2$), finding the target with $100\%$ certainty without any approximations!
  \item What happens to the success probability if you mistakenly run a second Grover iteration?
\end{enumerate}
\end{exercise}

\begin{exercise}\label{ch42:ex-no-cloning-overlap}
\textbf{(No-Cloning and inner product preservation.)}
An alternative proof of the No-Cloning Theorem proceeds by considering the inner product of two cloned states.
Suppose a unitary operator $U$ clones two distinct states $\ket{\phi}$ and $\ket{\psi}$:
\[
  U(\ket{\phi}\ket{0}) = \ket{\phi}\ket{\phi}, \qquad U(\ket{\psi}\ket{0}) = \ket{\psi}\ket{\psi}.
\]
\begin{enumerate}
  \item Using the fact that unitary operators preserve inner products, compute the inner product between the left-hand sides:
  $\braket{U(\phi\otimes 0)}{U(\psi\otimes 0)}$.
  \item Compute the inner product between the right-hand sides:
  $\braket{\phi\otimes\phi}{\psi\otimes\psi}$.
  \item Equate the two results to derive the equation $\braket{\phi}{\psi} = (\braket{\phi}{\psi})^2$.
  \item Conclude that cloning is only possible if $\braket{\phi}{\psi} = 0$ (orthogonal states) or $\abs{\braket{\phi}{\psi}} = 1$ (identical states), proving that an unknown arbitrary quantum state cannot be cloned.
\end{enumerate}
\end{exercise}
